% Figure: a thermistor calibration. Two panels: resistance vs.
% temperature on linear axes (left) and on Arrhenius-style axes (ln
% R vs 1/T) with a straight-line fit (right). The slope of the fit
% gives the B-constant of the thermistor model.
\begin{figure}[htbp]
\centering
\resizebox{\linewidth}{!}{%
\begin{tikzpicture}[scale=1.0,
  ax/.style={->, draw=engblack, thick},
  tick/.style={draw=engblack!60, thin},
  curve/.style={draw=engred, thick},
  fitline/.style={draw=engblue, thick, dashed},
  pt/.style={circle, fill=engred, inner sep=1.2pt},
  ann/.style={font=\footnotesize, align=center},
]

% --- Left: linear axes, R vs T (in degrees C) ---
\begin{scope}[xshift=0cm]
  \draw[ax] (0, 0) -- (6.0, 0) node[below right, ann] {$T$ ($\degC$)};
  \draw[ax] (0, 0) -- (0, 4.2) node[left, ann] {$R$ (k$\Omega$)};
  \foreach \xc/\xl in {0/0, 1/10, 2/20, 3/30, 4/50, 5/80} {
    \draw[tick] (\xc, -0.05) -- (\xc, 0.05);
    \node[ann, below, yshift=-2pt] at (\xc, 0) {\xl};
  }
  \foreach \y/\yl in {1/10, 2/30, 3/100, 4/300} {
    \draw[tick] (-0.05, \y) -- (0.05, \y);
    \node[ann, left, xshift=-2pt] at (0, \y) {\yl};
  }
  % Data points (R, T): T=0, 10, 20, 25, 30, 50, 80
  % R = 10 k * exp(B*(1/T_K - 1/298))   with B = 3800 K
  % Plotted at scaled coordinates
  \foreach \x/\y in {0/3.8, 1/3.3, 2/2.7, 2.4/2.45, 3/2.05, 4/1.25, 5/0.65} {
    \node[pt] at (\x, \y) {};
  }
  \draw[curve, thick] plot[smooth] coordinates
    {(0, 3.8) (1, 3.3) (2, 2.7) (2.4, 2.45) (3, 2.05) (4, 1.25) (5, 0.65)};
  \node[ann, anchor=south] at (3.5, 3.6) {(a) $R$ vs $T$, linear};
\end{scope}

% --- Right: Arrhenius axes, ln(R/R_0) vs 1/T ---
\begin{scope}[xshift=8cm]
  \draw[ax] (0, 0) -- (5.5, 0)
    node[below right, ann] {$1000 / T$ (K$^{-1}$)};
  \draw[ax] (0, -1) -- (0, 4.0) node[left, ann] {$\ln(R/R_{25})$};
  \foreach \xc/\xl in {0.5/2.8, 1.5/3.0, 2.5/3.2, 3.5/3.4, 4.5/3.6} {
    \draw[tick] (\xc, -0.05) -- (\xc, 0.05);
    \node[ann, below, yshift=-2pt] at (\xc, 0) {\xl};
  }
  \foreach \y/\yl in {-1/-1, 0/0, 1/1, 2/2, 3/3} {
    \draw[tick] (-0.05, \y) -- (0.05, \y);
    \node[ann, left, xshift=-2pt] at (0, \y) {\yl};
  }
  % Linear in 1/T: B*(1/T - 1/T_ref) ~ slope * (x - x_ref)
  % Fit slope ~ 3.8 K, so ln(R/R_25) ~ slope*(1000/T - 1000/298)
  % Plot 7 points on a line of slope 3.8 with offsets
  \foreach \x/\y in {0.5/-0.5, 1.5/-0.1, 2.5/0.6, 3.0/1.0, 3.5/1.5, 4.5/2.5} {
    \node[pt] at (\x, \y) {};
  }
  \draw[fitline] (0.3, -0.65) -- (4.7, 2.65);
  \node[ann, anchor=west, draw=engblue, fill=white, inner sep=2pt]
    at (0.6, 3.2) {fit: $\ln(R/R_{25}) = B(1/T - 1/T_{25})$\\
    slope $\approx B \approx 3800$ K};
  \node[ann, anchor=south] at (3.5, 3.6) {(b) Arrhenius axes};
\end{scope}

\end{tikzpicture}%
}%
\caption[Thermistor calibration: linear vs Arrhenius axes]{A
negative-temperature-coefficient (NTC) thermistor's resistance falls
with temperature according to the Steinhart approximation
$R(T) = R_{25} \exp\bigl[B(1/T - 1/T_{25})\bigr]$, with $T$ in
kelvin, $R_{25}$ the resistance at $25\,\degC$, and $B$ a
material constant. On linear axes (a) the data is a smooth decreasing
curve from which $B$ is hard to read. On Arrhenius axes (b),
plotting $\ln(R/R_{25})$ against $1/T$, the same data is a straight
line whose slope is exactly $B$. The example here uses a typical
$10\,\si{\kilo\ohm}$ NTC bead at $25\,\degC$ with $B = 3800\,\si
{\kelvin}$, calibrated across $0\,\degC$ to $80\,\degC$. The
chapter's design exercise (Section~2.7) asks the reader to set up
this same plot; the figure is a worked instance of the answer.}
\label{fig:vol02:ch02:thermistor-fit}
\end{figure}
